\documentclass[11pt]{article}
\usepackage[a4paper, top=18mm, bottom=18mm, left=20mm, right=20mm]{geometry}
\usepackage[T2A]{fontenc}
\usepackage[utf8]{inputenc}
\usepackage[ukrainian]{babel}
\usepackage{graphicx}
\usepackage[export]{adjustbox}
\usepackage{amsmath}
\usepackage{amssymb}
\graphicspath{{/app/Quiz/Scripts/2023/ProgAll/15BinTree/}{/app/Quiz/Scripts/2020/Mod2020_4/BinTree/}{/app/Quiz/Scripts/2021/Mod2021_3/BinTree/}}
\setlength{\parindent}{0pt}
\setlength{\parskip}{4pt}
\pagestyle{empty}
\begin{document}
%begin
{\large\textbf{Контрольна робота}}

\textbf{Варіант \No\,1}\hfill \textit{ПІБ:}\,\underline{\hspace{60mm}}
\vspace{4mm}

%beginitem
1. 
Скільки 2017-елементних підмножин множини натуральних від 1 до 20172017 містять принаймні одне число, що менше 172003
%enditem

%beginitem
2. 

a,b,c=9,6,8
def h(a,b=7,c=7):
    print(a,b,c,end="")

h(3,0,c=2)
print(a,b,c)


%enditem

%beginitem
3. 

\begin{center}
	\adjustimage{max size={0.8\linewidth}{0.8\paperheight}}
	{../15BinTree/trees_exp/100.png}
\end{center}
Указати послідовність міток вершин дерева, зображеного на рис., що утвориться в результаті обходу дерева в оберненому порядку. Мітки записувати через пробіл.\par Алгоритм починає роботу з кореня (на рис. це верхня вершина).
%answer:100 оберненому

%enditem

%beginitem
4. %goodpath:Scripts/2018/Mod2018_1/01-good.txt
Відмітити все, що є однією правильною лексемою:
( 1 ) x^2

( 2 ) false

( 3 ) +>

( 4 ) 'c'

%enditem

%beginitem
5. 
Записати ВИРАЗ, значенням якого є множина, що складається з кубівцілих чисел від 12 до 2002 (включно), що не діляться на жодне з чисел 2, 3.
%enditem

%beginitem
6. Що буде виведено за виконання фрагменту коду?

%code
#include <iostream>

int h(int a, int b){
    int c = 95;
    if (b >= -3)
        c = 3;
    else if (a < -5)
         return 0;
    else 
        c = 9;
    return c;
}

int main(){
    std::cout << h(-7, -7);
    return 0;
}
%code

%enditem

%beginitem
7. Що буде виведено за виконання фрагменту коду?

%code
#include <iostream>
using namespace std;

bool f(int n){
    cout<<"f";
    return n>=-2;
}

int main(){
    cout<<(f(-7) || f(-9));
    return 0;
}
%code

%enditem

%beginitem
8. Що буде виведено за виконання фрагменту коду?

%code
print('R', end=' ')
n=7
while n>=3:
    n+=-2
print(n)
%code


%enditem

%beginitem
9. Що буде виведено за виконання фрагменту коду?

%code
cout << 5 % 5 % 5 % 5;
%code

%enditem

%beginitem
10. 
try:
    res = 5//5%5*5
    if isinstance(res, bool):
        print(f'bool;{res}')
    elif isinstance(res, int):
        print(f'int;{res}')
    elif isinstance(res, float):
        print(f'float;{res:.2f}')
    else:
        print('???')
except:
    print('Z')

%enditem

%beginitem
11. 
a = 3
b = 0
c = 9

def h(b):
global c    a = 3
    b = 1
    c = 4
    return a + b + c

a = 7
b = 8
c = 1

try:
    print(h(b), a, b, c, sep=';')
except:
    print('ERR', a, b, c, sep=';')

%enditem

%beginitem
12. Що буде виведено за виконання фрагменту коду?

%code
print("3j" >= 1)
%code

%enditem

%beginitem
13. Що буде виведено за виконання фрагменту коду?
code
__d = 0

class D:
    __d = 1
    
    def __init__(self):
        self.__d = 2
        __d = 3
        D.__d = 4

obj = D()
try:
    print(__d, end=' ')
    print(obj.__d, end=' ')
    print(D.__d, end=' ')
    print(obj.__d, end=' ')
    print(obj._D__d, end=' ')
except:
    print('ERR')
code

%enditem

%beginitem
14. Що буде виведено за виконання фрагменту коду?

%code
d = [0,1,2,3,4,5,6,7,8,9]
print(d[-13:-13:-3])
%code


%enditem

%beginitem
15. 
Рядки журналу через двокрапку містять ім'я користувача (ідентифікатор), час події,тип події, код події, наприклад
\begin{verbatim}
s="username : 2022-05-05 13:26 : ERROR : 404"
\end{verbatim}
у коді 
\begin{verbatim}
res = re.fullmatch('???', s) 
s = ''
code_str = ??? 
\end{verbatim}
чим треба замінити кожен з ???, щоб значенням code\_str був код події? 



%enditem

%beginitem
16. Що буде виведено за виконання фрагменту коду?

%code
print('R', end=' ')
d=('a','b','c','d','e',0,1,2,3)
print(d[-11])
%code


%enditem

%beginitem
17. Що буде виведено за виконання фрагменту коду?

%code
a,b,c=3,0,9
def h(b):
global c    a=3
    b=1
    c=4
    return a+b+c

a,b,c=7,8,1
print(h(b),a,b,c)
%code

%enditem

%beginitem
18. 
d=[0,1,2,3,4,5,6,7,8,9]; print(d[-13:-13:-3])
%enditem

%beginitem
19. %goodpath:Scripts/2019/Mod2019_1/Lex/03-good.txt
Відмітити все, що є коректним оператором С++:
( 1 ) y=/x

( 2 ) if (1>2): x=10;

( 3 ) "a"=6

( 4 ) f"]={x:5.2f}{"

%enditem

%beginitem
20. 
x,y,z=3,0,2;  x,x,z,y=9,z,x,6;  print(x,y,z)
%enditem

%beginitem
21. 
def f(n): print("f", end=""); return n>=-2
   print(f(-7) or f(-9))
%enditem

%beginitem
22. 
n=3
   while n<=7: n+=2
   print(n)
%enditem

%beginitem
23. Що буде виведено за виконання фрагменту коду?
code
d = 0
k = 1

class D:
    k = 2
    
    def __init__(self):
global k        self.k = 3
        k = 4
        D.k = 5

obj = D()
print(d, k, D.k, obj.k, sep=' ')
code

%enditem

%beginitem
24. Що буде виведено за виконання фрагменту коду?

%code
print('R', end=' ')
d = ('a','b','c','d','e',0,1,2,3)
print(d[-11])
%code


%enditem

%beginitem
25. 
Значенням змінної \kw{a} є цілочисловий масив типу $numpy.ndarray$ розмірів $5{\times}5$. На скільки збільшиться сума елементів масиву \kw{a} після виконання
\begin{verbatim}
a.T[0] += 3
\end{verbatim}


Якщо код некоректний, то в якості відповіді зазначити ERROR.



%enditem

%beginitem
26. 
$a$ є цілочисловим масивом типу $numpy.ndarray$ розмірів $5{\times}5$. Сума елементів масиву $a$ дорівнює 22. Чому дорівнює сума елементів масиву $a$ після виконання 
\begin{verbatim}
a -= 22
\end{verbatim}


Якщо код некоректний, то в якості відповіді зазначити ERROR.



%enditem

%beginitem
27. 
not5.0>=5 or not2<8 or False>2.0
%enditem

%beginitem
28. 
d=['0','1','2','3','4']; del d[-6:-6]; print(d)
%enditem

%beginitem
29. 
Змінні \kw{next} та \kw{prev} вузла списку позначають наступний та попередній за ним вузол відповідно. Починаючи з голови, у список записано послідовні цілі числа від~$-36$ до~$49$. Нехай змінна \kw{p} позначає вузол, що зберігає значення~$-7$.
Яке число зберігається у вузлі \kw{p.next.next.next.next.next} ?

%enditem

%beginitem
30. 
def f(arg, n):
    n = 3
    tmp = arg
    del tmp[-6:-6]
    return len(tmp)>3

try:
    d = ['0','1','2','3','4']
    n = 0
    f(d, n)
    print(n, end=';')
    for el in d:
        print(el, end=';')
except:
    print('Z')

%enditem

%beginitem
31. 
d=['0','1','2','3','4']; print(d.pop(-2),end=' '); print(d)
%enditem

%beginitem
32. Що буде виведено за виконання фрагменту коду?
%code

x,y,z=3,0,2
x,x,z,y=9,z,x,6
print(x,y,z)
%code

%enditem

%beginitem
33. Що буде виведено за виконання фрагменту коду?

%code
print(not5.0>=5 or not2<8 or False>2.0)
%code

%enditem

%beginitem
34. 

print('R', end=' ')
d=['0','1','2','3','4']
del d[-6:-6]
print(d)


%enditem

%beginitem
35. 
def f(arg, n):
    n = 3
    tmp = arg
    del tmp[-6:-6]
    return len(tmp)>3

try:
    d = ['0','1','2','3','4']
    n = 0
    f(d, n)
    print(n, end=';')
    for el in d:
        print(el, end=';')
except:
    print('ERR')

%enditem

%beginitem
36. Що буде виведено за виконання фрагменту коду?

%code
print("3j">=1)
%code

%enditem

%beginitem
37. Що буде виведено за виконання фрагменту коду?

%code
#include <iostream>
using namespace std;

int h(int a, int b){
    int c = 95;
    if (b >= -3)
        c = 3;
    else if (a < -5)
         return 0;
    else 
        c = 9;
    return c;
}

int main(){
    cout << h(-7, -7);
    return 0;
}
%code

%enditem

%beginitem
38. 

\begin{center}
	\adjustimage{max size={0.8\linewidth}{0.8\paperheight}}
	{../BinTree/trees_exp/100.png}
\end{center}
Указати послідовність міток вершин дерева, зображеного на рис., що утвориться в результаті обходу дерева в оберненому порядку. Мітки записувати через пробіл.\par Обхід дерева починається з кореня (на рис. це верхня вершина).
%answer:100 оберненому

%enditem

%beginitem
39. Що буде виведено за виконання фрагменту коду?
%code
#include <iostream>
using namespace std;

int f(int &x, int &y){
    x = 3;
    y-= 1;
    return y;
}

int main(){
    int a = 8, b = 2;
    b = f(b, b);
    cout << a << ":" << b <<':';
    {
        int a = 6, b = 5;
        cout << ((b>=5) || ((a-=2) >= 5)) << ':';
        cout << a << ":" << b << ':';
    }
    {
        inta = 9;
        cout << a << ':';
    }
    cout << a << ":" << b << endl;
    return 0;
}
%code


%enditem

%beginitem
40. 
a,b,c=3,0,9
    def h(b):
      global c      a=3
      b=1
      c=4
      return a+b+c
    a,b,c=7,8,1
    print(h(b),a,b,c)
%enditem

%beginitem
41. 
В умовах попередньої задачі було виконано p->next=p->next->next;

Чому дорівнює значення p->next->next->next->next->n ?
%enditem

%beginitem
42. Що буде виведено за виконання фрагменту коду?

%code
d=['0','1','2','3']
print(d.pop(-3), end=' ')
print(d)
%code


%enditem

%beginitem
43. Що буде виведено за виконання фрагменту коду?

%code
5//5%5*5
%code

%enditem

%beginitem
44. 

print('R', end=' ')
n=7
while n>=3:
    n+=-2
print(n)


%enditem

%beginitem
45. Що буде виведено за виконання фрагменту коду?

%code
print('R', end=' ')
d=['0','1','2','3','4']
del d[-6:-6]
print(d)
%code


%enditem

%beginitem
46. Що буде виведено за виконання фрагменту коду?

%code
a = 3
b = 0
c = 9

def h(b):
global c    a*=3
    b=1
    c=4
    return a+b+c

a, b, c = 7, 8, 1
try:
    print(h(b), a, b, c, sep=';')
except:
    print('Z', a, b, c, sep=';')
%code

%enditem

%beginitem
47. Що буде виведено за виконання фрагменту коду?

%code
#include <iostream>
using namespace std;

int a = 3, b = 0, c = 9;

int h(int &b){
int c;    a = 3;
    b += 1;
    c = 4;
    return a + b + c;
}

int main(){
    inta = 7;
    intb = 8;
    intc = 1;
    cout << h(b) << ':';
    cout << a << ':' << b << ':' << c;
    return 0; 
}
%code

%enditem

%beginitem
48. 
a,b,c=3,0,9
    def h(b):
      global c      a*=3
      b=1
      c=4
      return a+b+c
    a,b,c=7,8,1
    print(h(b),a,b,c)
%enditem

%beginitem
49. Що буде виведено за виконання фрагменту коду?

%code
try:
    n = 7
    while n>=3:
        n += -2
    print(n, end=';')
except:
    print('Z')

%code

%enditem

%beginitem
50. 

\begin{center}
	\adjustimage{max size={0.8\linewidth}{0.8\paperheight}}
	{../BinTree/trees_exp/100.png}
\end{center}
Указати послідовність міток вершин дерева, зображеного на рис., що утвориться в результаті обходу дерева в оберненому порядку. Мітки записувати через пробіл.\par Алгоритм починає роботу з кореня (на рис. це верхня вершина).
%answer:100 оберненому

%enditem

%beginitem
51. 
try:
    n = 7
    while n>=3:
        n += -2
    print(n, end=';')
except:
    print('Z')

%enditem

%beginitem
52. Що буде виведено за виконання фрагменту коду?

%code
def h(a,b):
    c=95
    if b>=-3:
        c=3
    elif a<-5:
         return 0
    else: 
        c=9
    return c

print(h(-7,-7))
%code

%enditem

%beginitem
53. Що буде виведено за виконання фрагменту коду?

%code
#include <iostream>
using namespace std;

int h(int d){
    int x = 95;
    if (d >= -3) 
        x = 3;
    else if (d < -5)
         return 0;
    else
         x = 9;
    return x;
}

int main(){
    cout << h(-7);
    return 0;
}
%code


%enditem

%beginitem
54. %goodpath:Scripts/2018/Mod2018_1/03-good.txt
Відмітити все, що є коректним оператором С++:
( 1 ) cin<<3;

( 2 ) return 'a'<'c';

( 3 ) y=3^^2;

( 4 ) cout<<"\n";

%enditem

%beginitem
55. Що буде виведено за виконання фрагменту коду?

%code
def h(d):
    x=95
    if d>=-3: 
        x=3
    elif d<-5:
         return 0
    else:
         x=9
    return x

print(h(-7))
%code

%enditem

%beginitem
56. Що буде виведено за виконання фрагменту коду?

%code
cout << (!5.0>=5 || !2< 8 || false>2.0);
%code

%enditem

%beginitem
57. Що буде виведено за виконання фрагменту коду?

%code
print('R', end=' ')
d = [0,1,2,3,4,5,6,7,8,9]
print(d[-13:-13:-3])
%code


%enditem

%beginitem
58. 

print('R', end=' ')
d=['0','1','2','3','4']
print(d.pop(-2), end=' ')
print(d)


%enditem

%beginitem
59. Що буде виведено за виконання фрагменту коду?

%code
print('R', end=' ')
d=['0','1','2','3','4']
print(d.pop(-2), end=' ')
print(d)
%code


%enditem

%beginitem
60. Що буде виведено за виконання фрагменту коду?

%code
print(5//5%5*5)
%code

%enditem

%beginitem
61. 
try:
    n = 7
    while n>=3:
        n += -2
    print(n, end=';')
except:
    print('ERR')

%enditem

%beginitem
62. 
try:
    for d in range(-13, -13, -3):
        if d>=6:
            break
        print(d, end=';')
        if d>=6:
            break
    else:
        print(d,end=';')
    print(d)
except:
    print('ERR')

%enditem

%beginitem
63. 
"3j">=1
%enditem

%beginitem
64. Що буде виведено за виконання фрагменту коду?
%Оригинал был утерян. Требует отладки!
%code
__d = 0

class D:
    __d = 1
    
    def __init__(self):
        self.__d = 2
        __d = 3
        D.__d = 4

obj = D()
D.__d = 5
try:
    print(__d, end=' ')
    print(obj.__d, end=' ')
    print(D.__d, end=' ')
    print(obj.__d, end=' ')
    print(obj._D__d)
except:
    print('ERR')
%code

%enditem

%beginitem
65. 
У папці X31 розташовано проект, до якого входить синаксично правильна одиниця трансляції 1.cpp з рядком\par
{\#}include "2.h"
\parФайла 2.h у папці X31 нема, але він є в папці YZ. Що треба зробити, щоб одиниця трансляції 1.cpp, могла успішно відкомпілюватися?Переміщувати, копіювати та змінювати файли з кодом заборонено.\par\vspace{1.5cm}
%enditem

%beginitem
66. 

print('R', end=' ')
for d in range(-13,-13,-3):
    if d>=6:
        break
        print(d, end=' ')
    if d>=6:
        break
    else:
        print(d, end=' ')
    print(d)

%enditem

%beginitem
67. Що буде виведено за виконання фрагменту коду?

%code
n=3
while n<=7:
    n+=2
print(n, end=' ')
%code


%enditem

%beginitem
68. 
Записати ВИРАЗ, значенням якого є словник, ключами якого є цілі числа від 12 до 2002 (включно), що не діляться на жодне з чисел 2, 3, а ключам відповідають їх квадратні корені 
%enditem

%beginitem
69. 
Написати програму, що запитує у користувача значення дійсних параметрів $x$ та $y$, обчислює для них значення виразу 
$$\frac{\log_{2} x - 49}{(\tg x - 0.3) (y - 46)}$$ 
та повiдомляє введенi данi та результат обчислення.
 Оцінюється правильність та якість коду.


%enditem

%beginitem
70. %goodpath:Scripts/2019/Mod2019_1/Lex/01-good.txt
Відмітити все, що є однією правильною лексемою:
( 1 ) ?qwerty

( 2 ) z13

( 3 ) +>

( 4 ) %

%enditem

%beginitem
71. 
Змінні $next$ та $prev$ вузла списку позначають наступний та попередній за ним вузол відповідно. Починаючи з голови, у список записано послідовні цілі числа від $-36$ до $49$. Нехай змінна $p$ позначає вузол, що зберігає значення $-7$.
Яке значення зберігається у вузлі $p.next.next.next.next.next$ ?

%enditem

%beginitem
72. 
for d in range(-13,-13,-3):
    if d>=6:
        break
        print(d,end=' ')
    if d>=6:
        break
    else:
        print(d,end=' ')
    print(d)

%enditem

%beginitem
73. 
(Задача за частиною Пашка А.О. Наводиться в авторській редакції.)\par {\hspace {0.5 cm}}³дбудеться деяка подія $W$ чи ні залежить від двох факторів $A$ та $X$. За результатами статистичного експерименту (100 раз для кожного значення фактора $A$,  200 раз для кожного значення фактора $X$) подія $W$ відбулася:\par {\hspace {0.5 cm}}$(n+17)$ раз для значення $A(a)$, $(n+21)$ раз для значення $A(b)$, $(n+50)$ раз для значення $A(c)$, $(n+57)$ раз для значення $A(d)$, та відповідно $(2n+100)$ раз для значення $X(x)$, $(4n+17)$ раз для значення $X(y)$, $(n+150)$ раз для значення $X(z)$.\par
{\hspace {0.5 cm}}Знайти ймовірність того що подія $W$ не відбулася (умовна ймовірність) при значеннях факторів $A(d)$, $X(x)$. Фактори є незалежними, $n$ – номер цього екзаменаційного білета.\par
%enditem

%beginitem
74. 
for d in range(-13,-13,-3):
    if d>=6:
        break
        print(d,end=' ')
    if d>=6:
        break
else:
    print(d, end=' ')
print(d, end=' ')

%enditem

%beginitem
75. 
Записати ВИРАЗ, значенням якого є список, що складається з кубівцілих чисел від 12 до 2002 (включно), що не діляться на жодне з чисел 2, 3, записаних в оберненому порядку.
%enditem

%beginitem
76. 

def h(d):
    x=95
    if d>=-3: 
        x=3
    elif d<-5:
         return 0
    else:
         x=9
    return x

print(h(-7))

%enditem

%beginitem
77. 
Записати ВИРАЗ, значенням якого є словник, ключами якого є цілі числа від 12 до 2002 (включно), що не діляться на жодне з чисел 2, 3, а ключам відповідають їх квадратні корені 
%enditem

%beginitem
78. 
n=7
   while n>=3: n+=-2
   print(n)
%enditem

%beginitem
79. Що буде виведено за виконання фрагменту коду?
%code
for d in range(-13,-13,-3):
    if d>=6:
        break
        print(d,end=' ')
    if d>=6:
        break
else:
    print(d, end=' ')
print(d, end=' ')
%code

%enditem

%beginitem
80. 
Проект складається з од���ї одиниці трансляції 1.cpp, яка починається таким фрагментом\par
long h(){ return my134();}\par
Функція my134(); означена в 2.cpp.\par
Що треба зробити для того, щоб 1.cpp успішно відкомпілювався?\par\vspace{1.5cm}
Що треба зробити для того, щоб за проектом зібрався виконуваний код?\par

%enditem

%beginitem
81. Що буде виведено за виконання фрагменту коду?

%code
cout << (!5.0 >= 5 || !2 < 8 || false > 2.0);
%code

%enditem

%beginitem
82. 
Записати ВИРАЗ, значенням якого є список, що складається з кубівцілих чисел від 12 до 2002 (включно), що не діляться на жодне з чисел 2, 3, записаних в оберненому порядку.
%enditem

%beginitem
83. Що буде виведено за виконання фрагменту коду?

%code
for d in range(-8,-8+4,3):
    if d>=6:
        pass
    print(d, end=' ')
    d=-8
    if d>=6:
        break
else:
    print(d, end=' ')
print(d, end=' ')
%code

%enditem

%beginitem
84. 
try:
    res = 4**2.0-2
    if isinstance(res, bool):
        print(f'bool;{res}')
    elif isinstance(res, int):
        print(f'int;{res}')
    elif isinstance(res, float):
        print(f'float;{res:.2f}')
    else:
        print('???')
except:
    print('ERR')

%enditem

%beginitem
85. Що буде виведено за виконання фрагменту коду?

%code
print(5.0>5>2>8)
%code

%enditem

%beginitem
86. Що буде виведено за виконання фрагменту коду?

%code
4**2.0-2
%code

%enditem

%beginitem
87. 
Знайти кількість відношень на n-елементній множині, які неантирефлексивні або неантисиметричні
%enditem

%beginitem
88. Що буде виведено за виконання фрагменту коду?
%code
if (3 >= 3)
    cout << "d";
else
    cout << "d";
%code


%enditem

%beginitem
89. Що буде виведено за виконання фрагменту коду?

%code
"3j">=1
%code

%enditem

%beginitem
90. 

a,b,c=3,0,9
def h(b):
global c    a*=3
    b=1
    c=4
    return a+b+c

a,b,c=7,8,1
print(h(b),a,b,c)

%enditem

%beginitem
91. Що буде виведено за виконання фрагменту коду?

%code
cout << (5.0 >= 5 < 2 > 8);
%code

%enditem

%beginitem
92. 
a = 3
b = 0
c = 9

def h(b):
global c    a = 3
    b = 1
    c = 4
    return a + b + c

a = 7
b = 8
c = 1

try:
    print(h(b), a, b, c, sep=';')
except:
    print('Z', a, b, c, sep=';')

%enditem

%beginitem
93. Що буде виведено за виконання фрагменту коду?

%code
print('R', end=' ')
n=3
while n<=7:
    n+=2
print(n)
%code


%enditem

%beginitem
94. 
try:
    res = not5.0>=5 or not2<8 or False>2.0
    if isinstance(res, bool):
        print(f'bool;{res}')
    elif isinstance(res, int):
        print(f'int;{res}')
    elif isinstance(res, float):
        print(f'float;{res:.2f}')
    else:
        print('???')
except:
    print('Z')

%enditem

%beginitem
95. 
try:
    res = "3j">=1
    if isinstance(res, bool):
        print(f'bool;{res}')
    elif isinstance(res, int):
        print(f'int;{res}')
    elif isinstance(res, float):
        print(f'float;{res:.2f}')
    else:
        print('???')
except:
    print('ERR')

%enditem

%beginitem
96. 
Написати програму, що запитує у користувача значення дійсних параметрів $x$ та $y$, обчислює для них значення виразу 
$$\frac{\log_{2} x - 49}{(\tg x - 0.3) (y - 46)}$$ 
та повiдомляє введенi данi та результат обчислення.


%enditem

%beginitem
97. Що буде виведено за виконання фрагменту коду?

%code
cout << 5 % 5 % 5 % 5;
%code

%enditem

%beginitem
98. 
a = 3
b = 0
c = 9

def h(b):
global c    a *= 3
    b = 1
    c = 4
    return a + b + c

a = 7
b = 8
c = 1

try:
    print(h(b), a, b, c, sep=';')
except:
    print('ERR', a, b, c, sep=';')

%enditem

%beginitem
99. Що буде виведено за виконання фрагменту коду?

%code
print('R', end=' ')
d=[0,1,2,3,4,5,6,7,8,9]
print(d[-13:-13:-3])
%code


%enditem

%beginitem
100. 

try:
    print(3,end="")
    print(0j>=3j,end="")
    print(9,end="")
except TypeError: print(7,end="")
except Exception: print(8,end="")
else: print(1,end="")
finally: print(2,end="")

%enditem

%end
\newpage
\end{document}

